\documentclass[12pt,a4paper]{article}

\usepackage[T2A]{fontenc}
\usepackage[utf8]{inputenc}
\usepackage[russian]{babel}
\usepackage{amsmath,amssymb,amsfonts}
\usepackage{booktabs}
\usepackage{longtable}
\usepackage{geometry}
\usepackage{array}
\usepackage{enumitem}
\usepackage{verbatim}

\geometry{margin=2.2cm}

\title{Diamond State Model: геометрическая модель интерпретации состояний}
\author{Проект анализа интегрального индекса состояния}
\date{\today}

\begin{document}
\maketitle

\section{Назначение документа}
Настоящий документ формализует идею геометрической интерпретации состояний в виде \textbf{Diamond State Model} (DSM). Эта модель предназначена не для замены самой формулы ИИС, а для построения \textbf{карты состояний}, на которой каждый сегмент данных рассматривается как точка или вектор в скрытом пространстве признаков.

Основная интуиция такова:
\begin{itemize}[leftmargin=1.25cm]
    \item данные задают вектор состояния;
    \item вектор откладывается от центра карты;
    \item направление вектора указывает на \textbf{грань алмаза};
    \item грань соответствует \textbf{макрорежиму};
    \item внутри одной грани могут существовать локальные кластеры и подрежимы.
\end{itemize}

В двумерной версии получается ромб и естественная карта из четырёх внешних секторов и центральной переходной зоны. В трёхмерной версии алмаз превращается в октаэдр, что даёт восемь граней и центральную зону, то есть девять теоретических областей.

\section{Обозначения}
Пусть:
\begin{itemize}[leftmargin=1.25cm]
    \item $I_t$ --- значение интегрального индекса состояния в момент времени или на сегменте $t$;
    \item $R_t$ --- ресурсность организма или вторая ортогональная ось карты;
    \item $D_t$ --- динамический компонент, если используется трёхмерная модель;
    \item $x_t$ --- вектор состояния;
    \item $c$ --- центр карты;
    \item $s_I, s_R, s_D$ --- масштабы нормировки по соответствующим осям;
    \item $\varepsilon > 0$ --- малый стабилизирующий параметр;
    \item $Q_{0.25}, Q_{0.75}$ --- нижний и верхний квартили;
    \item $\rho_t$ или $\omega_t$ --- радиальная мера выраженности состояния.
\end{itemize}

\section{Двумерная версия DSM}
\subsection{Вектор состояния}
В двумерной версии состояние определяется как
\begin{equation}
x_t=
\begin{pmatrix}
I_t\\
R_t
\end{pmatrix},
\qquad
I_t = IIS_{V5}(t),
\qquad
R_t = RES(t).
\end{equation}

Здесь ось $I_t$ отвечает за текущее интегральное состояние, а ось $R_t$ --- за запас регуляции или устойчивость системы.

\subsection{Центр карты}
Центр карты задаётся как эмпирический нейтральный центр:
\begin{equation}
c=
\begin{pmatrix}
c_I\\
c_R
\end{pmatrix},
\qquad
c_I=\operatorname{median}(I),
\qquad
c_R=\operatorname{median}(R).
\end{equation}

На практике вместо медианы можно использовать baseline-центр, если исследователь хочет жёстко привязать карту к состоянию покоя.

\subsection{Нормировка}
Чтобы одна ось не доминировала над другой, вводятся масштабы:
\begin{equation}
s_I = Q_{0.75}(I)-Q_{0.25}(I),
\qquad
s_R = Q_{0.75}(R)-Q_{0.25}(R).
\end{equation}

После этого строится нормированный и центрированный вектор:
\begin{equation}
u_t=
\begin{pmatrix}
u_{t,1}\\
u_{t,2}
\end{pmatrix}
=
\begin{pmatrix}
\dfrac{I_t-c_I}{s_I+\varepsilon}\\[8pt]
\dfrac{R_t-c_R}{s_R+\varepsilon}
\end{pmatrix}.
\end{equation}

Вектор $u_t$ показывает, куда и насколько далеко текущее состояние отклоняется от нейтральной точки.

\subsection{Алмаз как ромб}
Двумерный алмаз задаётся как множество
\begin{equation}
D=\left\{u\in\mathbb{R}^2:\ |u_1|+|u_2|\le 1\right\}.
\end{equation}

Это ромб в метрике $L_1$. Такая форма удобна тем, что естественно разделяет пространство на четыре направления относительно центра.

\subsection{Сила состояния}
Вводится радиальная мера выраженности:
\begin{equation}
\rho_t = |u_{t,1}|+|u_{t,2}|.
\end{equation}

Смысл величины $\rho_t$:
\begin{itemize}[leftmargin=1.25cm]
    \item если $\rho_t$ мала, точка расположена близко к центру, и состояние выражено слабо;
    \item если $\rho_t$ велика, состояние выражено сильнее и тяготеет к одной из внешних граней карты.
\end{itemize}

Если $\rho_t > 0$, то пересечение луча с границей алмаза задаётся как
\begin{equation}
\hat{u}_t=\frac{u_t}{\rho_t}.
\end{equation}

Именно эта точка отвечает интуиции \textit{``вектор попал в грань алмаза''}.

\subsection{Пятирежимная интерпретация}
В двумерной версии удобно ввести центральную переходную зону и четыре внешних режима:
\begin{equation}
S_t=
\begin{cases}
\text{transition}, & \rho_t < \tau_0,\\
\text{regulated}, & u_{t,1}\ge 0,\ u_{t,2}\ge 0,\ \rho_t\ge \tau_0,\\
\text{mobilized}, & u_{t,1}<0,\ u_{t,2}\ge 0,\ \rho_t\ge \tau_0,\\
\text{depleted}, & u_{t,1}<0,\ u_{t,2}<0,\ \rho_t\ge \tau_0,\\
\text{fragile}, & u_{t,1}\ge 0,\ u_{t,2}<0,\ \rho_t\ge \tau_0.
\end{cases}
\end{equation}

Здесь $\tau_0$ --- порог центральной зоны. Его можно выбирать:
\begin{itemize}[leftmargin=1.25cm]
    \item фиксированно;
    \item как квантиль по распределению $\rho_t$;
    \item как порог, найденный по устойчивой кластерной структуре.
\end{itemize}

\subsection{Интерпретация секторов}
\begin{longtable}{>{\raggedright\arraybackslash}p{3.1cm} >{\raggedright\arraybackslash}p{10.5cm}}
\toprule
Режим & Интерпретация \\
\midrule
\endfirsthead
\toprule
Режим & Интерпретация \\
\midrule
\endhead
transition & Центральная зона, где состояние близко к нейтральному и выражено слабо. \\
regulated & Высокий ИИС и высокий ресурс. Система одновременно собрана и устойчива. \\
mobilized & ИИС снижен, но ресурс остаётся высоким. Возможна мобилизация под нагрузкой без явного срыва. \\
depleted & ИИС снижен и ресурс снижен. Наиболее напряжённая или истощённая область карты. \\
fragile & ИИС относительно высок, но ресурс низок. Внешне состояние может выглядеть приемлемым, но устойчивость системы слабая. \\
\bottomrule
\end{longtable}

\subsection{Мера уверенности}
Чтобы оценить, насколько уверенно точка принадлежит выбранному режиму, можно использовать:
\begin{equation}
\kappa_t = \rho_t - \tau_0.
\end{equation}

Чем больше $\kappa_t$, тем дальше точка от центральной переходной зоны и тем надёжнее назначение режима.

Дополнительно можно вводить коэффициент сбалансированности:
\begin{equation}
\beta_t=
\frac{\min(|u_{t,1}|,|u_{t,2}|)}
{\max(|u_{t,1}|,|u_{t,2}|)+\varepsilon}.
\end{equation}

Смысл:
\begin{itemize}[leftmargin=1.25cm]
    \item $\beta_t \approx 1$ означает, что точка идёт вдоль диагонали и режим выражен согласованно по двум осям;
    \item $\beta_t \approx 0$ означает, что точка почти лежит на одной оси и состояние более смешанное.
\end{itemize}

\subsection{Сетка и кластеры}
Если нужна именно сетка, а не только ромб, можно разбить пространство на ячейки:
\begin{equation}
G_{ij}=
\left\{
u\in D:
i\Delta_x \le u_1 < (i+1)\Delta_x,\ 
j\Delta_y \le u_2 < (j+1)\Delta_y
\right\}.
\end{equation}

Число наблюдений в ячейке:
\begin{equation}
n_{ij} = \#\{t:\ u_t\in G_{ij}\}.
\end{equation}

Доля режима $m$ в ячейке:
\begin{equation}
p_{ij}(m)=
\frac{\#\{t:\ u_t\in G_{ij},\ S_t=m\}}
{n_{ij}}.
\end{equation}

Таким образом:
\begin{itemize}[leftmargin=1.25cm]
    \item ромб задаёт макрорежимы;
    \item сетка показывает локальную плотность;
    \item кластеризация внутри сектора позволяет искать подрежимы.
\end{itemize}

\subsection{Схема 2D-карты}
\begin{verbatim}
                 regulated
                     /\
                    /  \
                   /    \
         mobilized /  +   \ fragile
                  /        \
                  \        /
                   \      /
                    \    /
                     \  /
                      \/
                   depleted

 + = transition core
\end{verbatim}

\section{Трёхмерная версия DSM}
\subsection{Третья ось}
Чтобы перейти от пяти зон к девяти теоретическим областям, вводится третья ось:
\begin{equation}
x_t=
\begin{pmatrix}
I_t\\
R_t\\
D_t
\end{pmatrix}.
\end{equation}

Наиболее естественный кандидат для $D_t$ --- динамический компонент:
\begin{equation}
D_t=
\tanh\!\left(
\lambda_1 \frac{dI}{dt}
\lambda_2 \frac{dR}{dt}
\lambda_3 \frac{d^2I}{dt^2}
\lambda_4 \,\mathrm{Vol}_t
\right),
\end{equation}
где:
\begin{itemize}[leftmargin=1.25cm]
    \item $\frac{dI}{dt}$ --- скорость изменения ИИС;
    \item $\frac{dR}{dt}$ --- скорость изменения ресурса;
    \item $\frac{d^2I}{dt^2}$ --- ускорение или кривизна траектории;
    \item $\mathrm{Vol}_t$ --- локальная волатильность.
\end{itemize}

В практической реализации не обязательно использовать все четыре члена одновременно; достаточно устойчивой и воспроизводимой динамической оси.

\subsection{Нормированный вектор}
Вводятся центр и масштабы:
\begin{equation}
c=
\begin{pmatrix}
c_I\\
c_R\\
c_D
\end{pmatrix},
\qquad
s_D=Q_{0.75}(D)-Q_{0.25}(D).
\end{equation}

После этого:
\begin{equation}
v_t=
\begin{pmatrix}
\dfrac{I_t-c_I}{s_I+\varepsilon}\\[8pt]
\dfrac{R_t-c_R}{s_R+\varepsilon}\\[8pt]
\dfrac{D_t-c_D}{s_D+\varepsilon}
\end{pmatrix}.
\end{equation}

\subsection{Алмаз как октаэдр}
Трёхмерный алмаз определяется как
\begin{equation}
O=
\left\{
v\in\mathbb{R}^3:
|v_1|+|v_2|+|v_3|\le 1
\right\}.
\end{equation}

Это правильный октаэдр в метрике $L_1$.

Радиальная мера:
\begin{equation}
\omega_t = |v_{t,1}|+|v_{t,2}|+|v_{t,3}|.
\end{equation}

\subsection{Девять теоретических областей}
Если $\omega_t < \tau_0$, то точка относится к центральному состоянию. Если $\omega_t \ge \tau_0$, то макрорежим определяется знаком трёх координат:
\begin{equation}
\operatorname{sgn}(v_t)=
\left(
\operatorname{sgn}(v_{t,1}),
\operatorname{sgn}(v_{t,2}),
\operatorname{sgn}(v_{t,3})
\right).
\end{equation}

Знаковых комбинаций ровно восемь:
\begin{equation}
(+,+,+),\ (+,+,-),\ (+,-,+),\ (+,-,-),\ (-,+,+),\ (-,+,-),\ (-,-,+),\ (-,-,-).
\end{equation}

Следовательно, полная теоретическая карта имеет:
\begin{equation}
8\ \text{внешних граней} + 1\ \text{центр} = 9\ \text{областей}.
\end{equation}

Именно так идея \textit{``девяти состояний''} получает строгую геометрическую форму.

\subsection{Смысл девяти областей}
В общем виде:
\begin{itemize}[leftmargin=1.25cm]
    \item первая координата отвечает за знак и выраженность интегрального состояния;
    \item вторая --- за ресурс и устойчивость;
    \item третья --- за направление динамики: рост, спад, переход, нестабильность.
\end{itemize}

Поэтому восемь внешних граней можно трактовать как восемь ориентированных макрорежимов, а центр --- как нейтрально-переходную область.

\section{Что является доказательством, а что нет}
Diamond State Model сама по себе не доказывает существование фиксированного числа реальных физиологических состояний. Она даёт \textbf{геометрическую карту интерпретации}. Чтобы область карты считать подтверждённым режимом, нужны отдельные данные:
\begin{itemize}[leftmargin=1.25cm]
    \item заполненность этой области реальными точками;
    \item устойчивость области по субъектам и подвыборкам;
    \item статистическая различимость относительно соседних областей;
    \item воспроизводимость на независимых датасетах.
\end{itemize}

Иначе говоря:
\begin{itemize}[leftmargin=1.25cm]
    \item \textbf{теоретическая область} --- это часть геометрии модели;
    \item \textbf{подтверждённый режим} --- это область, занятая и поддержанная реальными данными.
\end{itemize}

\section{Практическая стратегия применения}
Для текущего проекта наиболее разумна следующая последовательность:
\begin{enumerate}[leftmargin=1.25cm]
    \item использовать двумерную версию DSM на плоскости $(IIS_{V5}, RES)$;
    \item сначала анализировать 5 макрозон: четыре сектора и центральную область;
    \item проверять, какие из этих зон реально заняты и статистически различимы;
    \item только после этого переходить к трёхмерной версии с динамической осью;
    \item лишь в трёхмерной версии обсуждать полноценную карту из 9 теоретических областей.
\end{enumerate}

\section{Итог}
Diamond State Model позволяет формализовать интуитивную идею \textit{``алмаза состояний''} следующим образом:
\begin{equation}
\text{данные}
\;\to\;
\text{вектор состояния}
\;\to\;
\text{проекция на алмаз}
\;\to\;
\text{грань = макрорежим}
\;\to\;
\text{кластер внутри грани = подрежим}.
\end{equation}

В двумерной версии эта схема естественно даёт пять интерпретационных зон. В трёхмерной версии она даёт девять теоретических областей, из которых восемь соответствуют внешним граням октаэдра, а одна --- центральной зоне.

Следовательно, идея алмаза может рассматриваться как \textbf{полноценная геометрическая модель интерпретации состояний}, а не просто как метафора.

\end{document}
